\section{Độ đo là giả thuyết}
\label{sec:ch06-do-do-la-gia-thuyet}

\index{phân phối nền}%
\index{baseline}%
Công thức tích phân đã đặt giả định mà các chương trước thường gặp dưới tên
khác vào đúng vị trí của nó: độ đo $\mu_{j,x}$. Nếu độ đo giữ tọa độ $j$ ở
$x_j$ rồi lấy trung bình các tọa độ còn lại theo phân phối dữ liệu, attribution
trở thành một kỳ vọng có điều kiện. Nếu nó giữ $x_j$ và lấy trung bình phần còn
lại theo phân phối biên, ta nhận một biến thể giả định độc lập. Một lựa chọn
khác cho biểu thức của partial dependence plot~\autocite{p17firstprinciples}.

Ba trường hợp không chỉ khác kỹ thuật lấy mẫu. Chúng gán trọng lượng cho các
đầu vào khác nhau, nên chúng trả lời những câu hỏi khác nhau về mô hình. Chẳng
hạn, phân phối biên có thể ghép giá trị của các đặc trưng theo cách dữ liệu
thật không có. Phân phối có điều kiện tránh một số tổ hợp như vậy, nhưng đưa
cấu trúc phụ thuộc của dữ liệu vào lời giải thích. Không có lựa chọn nào được
hợp thức hóa chỉ vì tích phân trong phương trình~\ref{eq:ch06-attribution-do-do}
được tính chính xác.

Điều này nối trực tiếp với chương~\ref{ch:shap} và chương~\ref{ch:attribution-theo-gradient}.
Phân phối nền của SHAP và baseline của Integrated Gradients đều là cách chọn
đối chiếu. Khung độ đo không xóa sự lựa chọn ấy. Nó cho một ngôn ngữ chung để
nói rằng lựa chọn nằm trong attribution, không ở bên ngoài attribution.
